% Copy only the figures you want into your Springer manuscript.
% These figure files are created by MAXQDA_LDA_GPT_Publication_Figures.ipynb.
% If your manuscript is two-column, consider using figure* for wide figures.

\begin{figure}[!htb]
\centering
\includegraphics[width=\linewidth]{figure_maxqda_lda_heatmap.pdf}
\caption{Document-level thematic correspondence between 15 MAXQDA codes and seven dominant LDA topics in the 231 retained Vietnamese reference-positive texts. Cells indicate the percentage of documents assigned each code within the corresponding LDA topic. Outlines indicate manually mapped conceptually related themes; categories may co-occur. This is a descriptive comparison, not a formal agreement score.}
\label{fig:maxqda_lda_heatmap}
\end{figure}

\begin{figure}[!htb]
\centering
\includegraphics[width=\linewidth]{figure_three_way_joint_display.pdf}
\caption{Mixed-methods joint display comparing dominant LDA topics, corresponding researcher-reviewed MAXQDA codes, and the most frequent GPT primary thematic annotations. MAXQDA proportions use document denominators; GPT frequencies use annotation-record denominators and must not be interpreted as independent text-level agreement scores.}
\label{fig:maxqda_lda_gpt_joint}
\end{figure}

% Use this as supplementary material, or after researcher approval of quotations
% and English translations, because seven bilingual excerpts need generous space.
\begin{figure}[!htb]
\centering
\includegraphics[width=\linewidth]{figure_bilingual_representative_quotations.pdf}
\caption{Illustrative Vietnamese excerpts taken verbatim from MAXQDA-coded passages, with provisional English translations requiring researcher verification. The seven examples are organized by dominant LDA topic and are not intended to represent population-level theme prevalence.}
\label{fig:maxqda_bilingual_excerpts}
\end{figure}
